\documentclass[10pt,a4paper]{amsart}
\usepackage{amsmath,amssymb,amsthm,graphicx}
% Retained preamble for the Stacks Project chapter snapshot in this directory.
% The chapter's own preamble.tex is a copy of the shared Stacks Project preamble; the block
% below is sliced MECHANICALLY, line for line, from it. Nothing was added, removed or
% redefined; the wrapper class, the xr-hyper/hyperref wiring and the externaldocument list
% are omitted because the assembler supplies the document class.
%
% lines taken from preamble.tex: 3, 5, 9, 11, 12, 156, 157, 158, ... (51 lines total)
%
% package coverage restored from the chapter's own preamble (2026-10-13): xy, xy-2cell

\usepackage[all]{xy}
\xyoption{2cell}
\UseAllTwocells

\newenvironment{reference}{\comment}{\endcomment}
\newenvironment{slogan}{\comment}{\endcomment}
\newenvironment{history}{\comment}{\endcomment}
\theoremstyle{plain}
\newtheorem{theorem}[subsection]{Theorem}
\newtheorem{proposition}[subsection]{Proposition}
\newtheorem{lemma}[subsection]{Lemma}
\theoremstyle{definition}
\newtheorem{definition}[subsection]{Definition}
\newtheorem{example}[subsection]{Example}
\newtheorem{exercise}[subsection]{Exercise}
\newtheorem{situation}[subsection]{Situation}
\theoremstyle{remark}
\newtheorem{remark}[subsection]{Remark}
\newtheorem{remarks}[subsection]{Remarks}
\numberwithin{equation}{subsection}
\def\lim{\mathop{\mathrm{lim}}\nolimits}
\def\colim{\mathop{\mathrm{colim}}\nolimits}
\def\Spec{\mathop{\mathrm{Spec}}}
\def\Hom{\mathop{\mathrm{Hom}}\nolimits}
\def\Ext{\mathop{\mathrm{Ext}}\nolimits}
\def\SheafHom{\mathop{\mathcal{H}\!\mathit{om}}\nolimits}
\def\SheafExt{\mathop{\mathcal{E}\!\mathit{xt}}\nolimits}
\def\Sch{\mathit{Sch}}
\def\Mor{\mathop{\mathrm{Mor}}\nolimits}
\def\Ob{\mathop{\mathrm{Ob}}\nolimits}
\def\Sh{\mathop{\mathit{Sh}}\nolimits}
\def\NL{\mathop{N\!L}\nolimits}
\def\CH{\mathop{\mathrm{CH}}\nolimits}
\def\proetale{{pro\text{-}\acute{e}tale}}
\def\etale{{\acute{e}tale}}
\def\QCoh{\mathit{QCoh}}
\def\Ker{\mathop{\mathrm{Ker}}}
\def\Im{\mathop{\mathrm{Im}}}
\def\Coker{\mathop{\mathrm{Coker}}}
\def\Coim{\mathop{\mathrm{Coim}}}
\DeclareMathSymbol{\boxtimes}{\mathbin}{AMSa}{"02}
\def\QCohstack{\mathcal{QC}\!\mathit{oh}}
\def\Cohstack{\mathcal{C}\!\mathit{oh}}
\def\Spacesstack{\mathcal{S}\!\mathit{paces}}
\def\Quotfunctor{\mathrm{Quot}}
\def\Hilbfunctor{\mathrm{Hilb}}
\def\Curvesstack{\mathcal{C}\!\mathit{urves}}
\def\Polarizedstack{\mathcal{P}\!\mathit{olarized}}
\def\Complexesstack{\mathcal{C}\!\mathit{omplexes}}
\def\Pic{\mathop{\mathrm{Pic}}\nolimits}
\def\Picardstack{\mathcal{P}\!\mathit{ic}}
\def\Picardfunctor{\mathrm{Pic}}
\def\Deformationcategory{\mathcal{D}\!\mathit{ef}}
\title{Stacks, Deformation Theory: lemma-canonical-class-algebra}
\author{The Stacks Project Authors}
\date{stacks-project:defos.tex:local-snapshot-20261012; GNU FDL 1.2 or later}
\begin{document}
\maketitle
\noindent\textit{Editorial note.} Statement span is lines 5134-5159 of sources/stacks-defos/pinned-original-source.tex; proof span is lines 5161-5244, the complete author proof (82 body lines).\par\medskip
\noindent\textit{Editorial note.} Navigation stubs: 0 cited-but-not-included labels are supplied as non-load-bearing pointers with their pinned tags where the snapshot tag index has them.\par\medskip
\noindent\textit{Editorial note.} extra\_packages is empty on purpose: the retained preamble already loads xy with options and a second declaration clashes in the assembler wrapper.\par\medskip
\noindent\textit{Editorial note.} Text is retained verbatim, including the author's own typos and placeholders.\par\medskip
\noindent\textit{Editorial note.} Queue index 8 of queue\_defos.json (pb=3263, stmt 5134-5159, proof 5161-5244).\par\medskip
\noindent\textit{Allowed standard prerequisites (editorial).}\begin{itemize}
\item Explicit prerequisite boundary. The cited inputs are prerequisites; the argument in the retained span is not replaced by any cited result.
\end{itemize}\medskip
\begin{lemma}
\label{lemma-canonical-class-algebra}
Let $R' \to R$ be a surjection of rings whose kernel is an ideal
$I$ of square zero. For every $K \in D^-(R)$ there is a canonical
map
$$
\omega(K) : K \longrightarrow K \otimes_R^\mathbf{L} I[2]
$$
in $D(R)$ with the following properties
\begin{enumerate}
\item $\omega(K) = 0$ if and only if there exists
$K' \in D(R')$ with $K' \otimes_{R'}^\mathbf{L} R = K$,
\item given $K \to L$ in $D^-(R)$ the diagram
$$
\xymatrix{
K \ar[d] \ar[rr]_-{\omega(K)} & &
K \otimes^\mathbf{L}_R I[2] \ar[d] \\
L \ar[rr]^-{\omega(L)} & &
L \otimes^\mathbf{L}_R I[2]
}
$$
commutes, and
\item formation of $\omega(K)$ is compatible with ring maps $R' \to S'$
(see proof for a precise statement).
\end{enumerate}
\end{lemma}

\begin{proof}
Choose a bounded above complex $K^\bullet$ of free $R$-modules representing
$K$. Then we can choose free $R'$-modules $(K')^n$ lifting $K^n$.
We can choose $R'$-module maps $(d')^n_K : (K')^n \to (K')^{n + 1}$
lifting the differentials $d^n_K : K^n \to K^{n + 1}$ of $K^\bullet$.
Although the compositions
$$
(d')^{n + 1}_K \circ (d')^n_K : (K')^n \to (K')^{n + 2}
$$
may not be zero, they do factor as
$$
(K')^n \to K^n \xrightarrow{\omega^n_K}
K^{n + 2} \otimes_R I = I(K')^{n + 2} \to (K')^{n + 2}
$$
because $d^{n + 1} \circ d^n = 0$.
A calculation shows that $\omega^n_K$ defines a map of complexes.
This map of complexes defines $\omega(K)$.

\medskip\noindent
Let us prove this construction is compatible with a map of complexes
$\alpha^\bullet : K^\bullet \to L^\bullet$ of bounded above free $R$-modules
and given choices of lifts $(K')^n, (L')^n, (d')^n_K, (d')^n_L$.
Namely, choose $(\alpha')^n : (K')^n \to (L')^n$ lifting the
components $\alpha^n : K^n \to L^n$. As before we get a
factorization
$$
(K')^n \to K^n \xrightarrow{h^n}
L^{n + 1} \otimes_R I = I(L')^{n + 1} \to (L')^{n + 2}
$$
of $(d')^n_L \circ (\alpha')^n - (\alpha')^{n + 1} \circ (d')_K^n$.
Then it is an pleasant calculation to show that
$$
\omega^n_L \circ \alpha^n =
(d_L^{n + 1} \otimes \text{id}_I) \circ h^n + h^{n + 1} \circ d_K^n +
(\alpha^{n + 2} \otimes \text{id}_I) \circ \omega^n_K
$$
This proves the commutativity of the diagram in (2) of the lemma
in this particular case. Using this for two different choices
of bounded above free complexes representing $K$, we find that
$\omega(K)$ is well defined! And of course (2) holds in general as well.

\medskip\noindent
If $K$ lifts to $K'$ in $D^-(R')$, then we can represent
$K'$ by a bounded above complex of free $R'$-modules
and we see immediately that $\omega(K) = 0$.
Conversely, going back to our choices $K^\bullet$,
$(K')^n$, $(d')^n_K$, if $\omega(K) = 0$, then we can find
$g^n : K^n \to K^{n + 1} \otimes_R I$ with
$$
\omega^n = (d_K^{n + 1} \otimes \text{id}_I) \circ g^n +
g^{n + 1} \circ d_K^n
$$
This means that with differentials
$(d')^n_K - g^n : (K')^n \to (K')^{n + 1}$
we obtain a complex of free $R'$-modules lifting $K^\bullet$.
This proves (1).

\medskip\noindent
Finally, part (3) means the following: Let $R' \to S'$ be a map of
rings. Set $S = S' \otimes_{R'} R$ and denote $J = IS' \subset S'$
the square zero kernel of $S' \to S$. Then given $K \in D^-(R)$
the statement is that we get a commutative diagram
$$
\xymatrix{
K \otimes_R^\mathbf{L} S \ar[d] \ar[rr]_-{\omega(K) \otimes \text{id}} & &
(K \otimes^\mathbf{L}_R I[2]) \otimes_R^\mathbf{L} S \ar[d] \\
K \otimes_R^\mathbf{L} S \ar[rr]^-{\omega(K \otimes_R^\mathbf{L} S)} & &
(K \otimes_R^\mathbf{L} S) \otimes^\mathbf{L}_S J[2]
}
$$
Here the right vertical arrow comes from
$$
(K \otimes^\mathbf{L}_R I[2]) \otimes_R^\mathbf{L} S =
(K \otimes_R^\mathbf{L} S) \otimes_S^\mathbf{L}
(I \otimes_R^\mathbf{L} S)[2] \longrightarrow
(K \otimes_R^\mathbf{L} S) \otimes_S^\mathbf{L} J[2]
$$
Choose $K^\bullet$, $(K')^n$, and $(d')^n_K$ as above.
Then we can use $K^\bullet \otimes_R S$, $(K')^n \otimes_{R'} S'$, and
$(d')^n_K \otimes \text{id}_{S'}$ for the construction of
$\omega(K \otimes_R^\mathbf{L} S)$.
With these choices commutativity
is immediately verified on the level of maps of complexes.
\end{proof}

\end{document}
